\chapter{Introduction}

\section{Statement of the Problem}
\label{sec:statement-problem}

The proliferation of free online learning resources---YouTube tutorials, blog posts, open-source documentation, and platforms like FreeCodeCamp and GeeksForGeeks---has democratized access to technical knowledge. However, this abundance has created a paradoxical challenge: fragmentation without structure. Ethiopian students and self-directed learners face a critical problem landscape characterised by six interconnected challenges.

\textbf{Challenge 1: The Tutorial Hell Phenomenon.} Learners navigate an overwhelming landscape of resources without a coherent learning path. A student seeking to learn web development might start with an HTML tutorial, then jump to a CSS playlist on YouTube, then read scattered JavaScript documentation, then discover React, without understanding how these components interconnect or progress logically. This non-sequenced, ad-hoc learning creates knowledge gaps, misaligned prerequisites, and wasted effort. Research on cognitive load theory~\cite{sweller1988} demonstrates that unstructured, fragmented information presentation significantly increases cognitive load, reducing learning efficiency and retention.

\textbf{Challenge 2: Absence of Adaptive Feedback Loops.} Existing free resources are predominantly passive. Learners consume content (watch videos, read documentation) but receive minimal performance-based feedback or adaptation. If a learner struggles with a concept, the resource does not diagnose why or adapt the explanation. Unlike human tutors who adapt in real-time based on learner confusion, static resources force learners to self-diagnose and manually seek alternative resources. This mirrors the broader finding in educational psychology that immediate, targeted feedback is essential for efficient learning~\cite{hattie2007}, yet most self-directed platforms lack this feature.

\textbf{Challenge 3: No Mastery-Based Progression Gates.} Learners can mark topics as \enquote{complete} based on time spent or arbitrary quizzes, regardless of actual competence. A learner might watch a JavaScript video, feel confident, and mark \enquote{JavaScript Fundamentals} as done---then struggle severely when prerequisites are required in an advanced topic. This leads to a false sense of mastery and downstream learning failures. Research on mastery-based progression~\cite{bloom1968} establishes that learning should be gated by demonstrated competence, not time or completion of content consumption.

\textbf{Challenge 4: Cognitive Overload and Retention Decay.} Self-directed learners often experience cognitive overload: learning multiple new concepts simultaneously without reinforcement, leading to rapid forgetting. Ebbinghaus's forgetting curve~\cite{ebbinghaus1885} demonstrates that without strategic review, learned information decays exponentially within days. Yet self-directed learners have no automated system prompting spaced repetition; they must manually re-engage with material, which is cognitively demanding and often neglected. Consequently, skills learned weeks or months prior decay to unusable levels by the time they are needed in integrated projects.

\textbf{Challenge 5: Generic Roadmaps Ignore Learner Heterogeneity.} Static, one-size-fits-all learning paths (e.g., \enquote{Learn HTML $\rightarrow$ CSS $\rightarrow$ JavaScript $\rightarrow$ React $\rightarrow$ Node.js}) ignore that learners differ in preferences, prior knowledge, and career aspirations. Some learners prefer visual, design-oriented paths; others prefer backend logic and databases. A fixed sequence frustrates learners whose interests diverge from the prescribed path, reducing intrinsic motivation and increasing dropout rates. Research on self-determination theory~\cite{ryan2000, deci1985} shows that learner agency---the ability to make meaningful choices about learning---significantly impacts motivation and persistence.

\textbf{Challenge 6: The Hallucination Problem in AI-Driven Recommendations.} Recent tools leveraging Large Language Models (LLMs) for curriculum generation can produce fluent but factually incorrect sequences. An LLM might generate a learning path for Python that contains logical inconsistencies, missing prerequisites, or inaccurate technical descriptions~\cite{marcus2020, weidinger2021}. In technical education, where precision is non-negotiable, hallucination in curriculum design is catastrophic. Learners might internalize incorrect concepts, leading to buggy code or flawed system design.

\textbf{Systemic Root Cause: Lack of Verifiable, Adaptive Knowledge Structures.} Underlying these six challenges is a fundamental structural problem: the absence of a verified, machine-readable knowledge graph that combines the rigor of traditional Intelligent Tutoring Systems with the scalability and flexibility of modern LLMs. Current solutions fall into two inadequate categories:
\begin{enumerate}
    \item Static resources (documentation, videos, blogs): high quality but passive, non-adaptive, and fragmented.
    \item Pure LLM systems (ChatGPT tutoring): flexible and scalable but prone to hallucination and lack curriculum coherence.
\end{enumerate}
Neither category addresses the integration of verified domain knowledge, adaptive progression, and performance-based feedback at the system level.

\section{Objectives}
\label{sec:objectives}

\subsection{General Objective}

To design, develop, and validate an ontology-guided, adaptive learning roadmap generator that integrates verified domain knowledge with intelligent tutoring principles to enable self-directed Ethiopian learners to achieve technical competence in programming and software engineering through structured, performance-driven, and personalized learning paths.

\subsection{Specific Objectives}

\textbf{Objective 1: Design and Implement a Hybrid AI Architecture.} Develop and operationalise the \enquote{Skeleton \& Flesh} hybrid architecture that separates curriculum knowledge representation (the verified ontology skeleton) from runtime tutoring (the adaptive AI flesh). Specifically:
\begin{itemize}
    \item Create domain-specific ontologies (DAG of learning nodes, prerequisites, learning outcomes) for five initial domains: Frontend Development (HTML/CSS/JavaScript/React), Backend Development (Node.js/Databases/APIs), Data Science (Python/Pandas/ML), DevOps Engineering (Docker/CI-CD/Cloud), and Python (core language fundamentals).
    \item Implement a Teacher Model (offline LLM-based generator) that produces the initial ontology with human expert verification.
    \item Implement a Tutor Model (runtime, three-tier AI stack) that generates explanations, adaptive quizzes, and conversational answers grounded in the ontology, preventing hallucination through skeleton-constrained generation.
\end{itemize}

\textbf{Objective 2: Implement the Gatekeeper Pattern for Assessment-Driven Progression.} Engineer and validate the Gatekeeper Pattern, an assessment-driven progression mechanism where learners cannot advance to subsequent learning nodes without passing performance thresholds on adaptive quizzes. Specifically:
\begin{itemize}
    \item Design three-tier quiz outcomes: Pass ($\ge 80\%$ $\rightarrow$ unlock next node + challenge project), Marginal Pass ($70$--$79\%$ $\rightarrow$ unlock node + flag for spaced repetition), Fail ($<70\%$ $\rightarrow$ trigger resource adaptation or prerequisite review).
    \item Implement adaptive content modality: if a learner fails ($<70\%$), trigger a \texttt{resource\_swap} adaptation event that instructs the Tutor Model to regenerate the node explanation in a more tutorial-oriented, step-by-step style rather than a reference-documentation style, recognising modality mismatch in learning.
    \item Implement prerequisite validation: if a learner fails ($<50\%$), recommend review of foundational nodes before retry.
\end{itemize}

\textbf{Objective 3: Integrate Spaced Repetition with Decay-Aware Roadmap Evolution.} Operationalise Mastery Decay---a system that integrates Ebbinghaus's forgetting curve directly into the learner's roadmap visualisation, making knowledge decay visible and triggering timely reviews. Specifically:
\begin{itemize}
    \item Implement a decay algorithm where nodes transition from \enquote{Mastered} (green, 0--14 days) to \enquote{Review Needed} (yellow, 14--30 days) to \enquote{Relearn} (red, $>30$ days) based on time since last visit.
    \item Generate micro-quizzes (2--3 questions) when nodes transition to yellow or red status, targeting core concepts of that node.
    \item Record and visualise decay patterns, enabling learners to identify their personal retention curves and adjust study frequency.
\end{itemize}

\textbf{Objective 4: Develop Multi-Path Branching for Personalized Learning.} Design and implement Multi-Path Branching where learners choose between alternative learning sequences based on interests and strengths, while ensuring semantic coherence and prerequisite integrity. Specifically:
\begin{itemize}
    \item After foundational nodes, offer learners a choice: \enquote{Prefer visual design or backend logic?} recommending either Frontend (CSS $\rightarrow$ UI frameworks) or Backend (Databases $\rightarrow$ APIs) paths.
    \item Ensure paths reconverge at advanced, integrative topics (e.g., \enquote{Full-Stack Architecture}, \enquote{System Design}).
    \item Validate that paths maintain prerequisite dependencies and learning outcome alignment.
\end{itemize}

\textbf{Objective 5: Develop AI-Grounded Resource Generation.} Build a quality-assured content generation pipeline that produces node-specific learning resources and references anchored entirely within the verified ontology skeleton. Specifically:
\begin{itemize}
    \item Generate node explanations and supplementary resource recommendations at runtime using the Tutor Model, constrained to the node's verified learning outcomes so that no external hallucinated links are introduced.
    \item Attach curated reference links to ontology nodes during the offline authoring phase, verified by domain experts before publication.
    \item Enable domain expert instructors to review and update node resources through the administrator ontology builder interface, maintaining content quality without system downtime.
\end{itemize}

\textbf{Objective 6: Develop Role-Stratified User Interfaces.} Create intuitive, responsive interfaces for all three system roles---learner, domain expert instructor, and administrator---enabling end-to-end interaction with the platform. Specifically:
\begin{itemize}
    \item Develop a full-featured web application (React~18 + TypeScript) providing learners with an interactive roadmap DAG, progress dashboards, quiz flows, streaming AI explanations, an AI Instructor chatbot, learning insights analytics, achievement tracking, and certificate download/sharing.
    \item Develop an instructor portal enabling domain experts to monitor learner progress, view domain-level analytics, review flagged quiz events, and manage ontology content.
    \item Develop an administrative console for user management, domain lifecycle management, ontology construction, and system statistics monitoring.
    \item Develop a Flutter-based mobile application (Android/iOS) as a complementary access channel, providing roadmap navigation and quiz interactions for learners on mobile devices.
\end{itemize}

\textbf{Objective 7: Validate System Effectiveness with Ethiopian Learners.} Conduct empirical evaluation of \projectname's learning efficacy and usability with real learners. Specifically:
\begin{itemize}
    \item Recruit 50--100 Ethiopian students/engineers as test users.
    \item Measure learning gains: assess pre/post knowledge using domain-aligned assessments, comparing learners using \projectname{} vs.\ a control group using traditional resources.
    \item Measure engagement and retention: track time-on-task, completion rates, session frequency, and long-term return visits (6+ months).
    \item Collect usability feedback via surveys and interviews, identifying UI/UX improvements and cultural/linguistic adaptations needed.
\end{itemize}

\textbf{Objective 8: Ensure System Scalability and Sustainability.} Design and implement \projectname{} as a scalable, maintainable, open-source system capable of extension to additional domains and deployment in resource-constrained environments. Specifically:
\begin{itemize}
    \item Deploy using Docker Compose with an Nginx reverse proxy, enabling reproducible, infrastructure-agnostic deployment on any VPS or cloud host without vendor lock-in.
    \item Implement circuit breakers on all three AI providers to guarantee graceful degradation and uninterrupted service when individual AI backends are unavailable.
    \item Maintain Redis-based response caching to minimise redundant AI API calls and reduce operating costs at scale.
    \item Create comprehensive documentation and developer guides enabling future contributors to extend the system to new domains or platforms.
\end{itemize}

\section{Scope and Limitations}
\label{sec:scope-limitations}

\subsection{In-Scope Elements}

The \projectname{} system encompasses the following capabilities and features:

\textbf{1. Learning Domains:} Five technical domains are in scope: (a) Frontend Development (HTML, CSS, JavaScript ES6+, React.js, responsive design), (b) Backend Development (Node.js, Express, RESTful API design, SQL/NoSQL databases), (c) Data Science (Python, Pandas, NumPy, basic ML algorithms, data visualisation), (d) DevOps Engineering (Docker, CI/CD pipelines, cloud infrastructure fundamentals), and (e) Python (core language fundamentals, OOP, standard library). These domains were selected as high-demand skills with well-defined prerequisite structures and representational breadth across modern software engineering.

\textbf{2. Hybrid AI Architecture:} Full implementation of the Skeleton \& Flesh strategy, including (a) a Teacher Model that uses LLMs to generate the initial ontology DAG with expert verification, (b) a runtime Tutor Model implementing a three-tier provider chain (Phi-4-Multimodal $\to$ Ollama/Qwen2.5 $\to$ Gemini 2.5 Flash) that generates explanations, adaptive quizzes, and conversational answers constrained by the verified ontology, and (c) circuit-breaker-guarded provider failover to ensure continuous availability when individual AI backends are unreachable.

\textbf{3. Gatekeeper Pattern:} Complete implementation of assessment-driven progression, including (a) adaptive quiz generation for each node, (b) three-tier outcome logic (Pass/Marginal Pass/Fail), (c) adaptive resource swapping triggered by failure modes, (d) prerequisite validation and enforcement, and (e) challenge project recommendations for learners who pass with high scores.

\textbf{4. Spaced Repetition Integration:} Full implementation of Mastery Decay, including (a) decay-state transitions (Green $\rightarrow$ Yellow $\rightarrow$ Red) based on time since last visit, (b) automated micro-quiz triggers, (c) visual roadmap representation of decay states, and (d) learner notifications prompting review.

\textbf{5. Multi-Path Branching:} Implementation of learner choice at key curriculum junctures, including (a) branching questions (e.g., \enquote{Prefer frontend or backend?}), (b) alternative learning sequences for each path, (c) prerequisites maintained within each path, and (d) reconvergence points at advanced topics.

\textbf{6. AI-Grounded Content and Resource Generation:} Full implementation of ontology-constrained content generation, including (a) streaming AI explanations delivered token-by-token via Server-Sent Events (SSE), (b) a conversational AI Instructor chatbot supporting real-time question-and-answer within each learning node, and (c) node-level resource references curated offline by domain experts and surfaced at runtime.

\textbf{7. User Interfaces:} Complete development of (a) a web application (React 18 + TypeScript) for desktop/laptop learners with roadmap DAG visualisation, progress dashboards, quiz flows, streaming AI panels, learning insights analytics, gamification achievements, and certificate management; (b) an instructor portal for learner monitoring, domain analytics, and flagged event review; (c) an admin console for user, domain, and ontology management; and (d) a Flutter mobile application (Android/iOS) as a complementary access channel.

\textbf{7a. Certificates:} Full implementation of verifiable completion certificates, including (a) server-issued certificates with Crockford base-32 public identifiers upon roadmap mastery, (b) client-side PDF and PNG download via html-to-image and jsPDF, and (c) a publicly accessible, no-authentication verification page at \texttt{/verify/:publicId} for sharing and employer validation.

\textbf{8. User Base:} Initial deployment targets Ethiopian software engineering students (university-level), self-taught engineers, and career-transitioning professionals. A target user base of 50--100 test users for initial validation, with potential to scale to thousands.

\textbf{9. Deployment and Scalability:} Deployment using Docker Compose orchestration with Nginx as a reverse proxy, providing infrastructure-agnostic containerized deployment suitable for any Linux VPS or cloud host. The microservices architecture (API Gateway, Learning Service, AI Service) with Redis caching enables horizontal scaling of individual services independently.

\textbf{10. Methodology:} Agile-Scrum development methodology with 2-week sprints, daily standups, sprint reviews, and retrospectives, ensuring iterative delivery and stakeholder feedback integration.

\subsection{Out-of-Scope Elements}

The following elements are explicitly excluded from this project:
\begin{enumerate}
    \item \textbf{Spoken Language Learning:} The system does not support learning of spoken languages. It focuses exclusively on programming languages and technical content, where precision and verification are more tractable.
    \item \textbf{Paid Content Integration:} The system does not aggregate links to paid courses or license proprietary content. It restricts to free, openly accessible resources, ensuring accessibility for Ethiopian learners with limited financial means.
    \item \textbf{Video Content Scraping:} The system does not scrape, host, or re-distribute video content. While it may link to high-quality YouTube videos from trusted educational channels, it does not download or modify video content.
    \item \textbf{Advanced Gamification:} While the system implements basic achievement tracking and progress visualisation, advanced competitive gamification elements (leaderboards, peer challenges, community rankings) are out of scope.
    \item \textbf{Institutional Accreditation:} The system issues verifiable digital certificates of completion with publicly accessible verification links. However, these certificates are not formally accredited by academic institutions or professional certification bodies in the initial deployment.
    \item \textbf{Collaborative Learning:} The system does not include peer collaboration features (discussion forums, pair programming rooms, group projects). Learning paths are individualised.
    \item \textbf{LMS Integration:} Integration with existing Learning Management Systems (Moodle, Canvas, Blackboard) is out of scope. The system operates as a standalone platform.
    \item \textbf{Advanced Personalisation Based on Learning Styles:} While Multi-Path Branching provides some personalisation, the system does not implement advanced learning style detection~\cite{pashler2008}.
\end{enumerate}

\subsection{Limitations}

The project acknowledges the following realistic constraints:
\begin{enumerate}
    \item \textbf{Ontology Verification Bottleneck:} The quality of the skeleton ontology depends on expert verification. If domain experts are unavailable or provide incomplete feedback, the ontology may contain errors or omissions.
    \item \textbf{Resource Variability:} The quality of AI-generated explanations and supplementary resources varies with LLM provider quality. If the primary and fallback AI providers are both unavailable, the system falls back to cached responses or static node content, which may be less tailored to the learner.
    \item \textbf{Algorithm Calibration:} The Mastery Decay algorithm parameters (14-day yellow threshold, 30-day red threshold, 80\% pass score) are initial estimates informed by research but may require calibration based on real learner data.
    \item \textbf{Scalability Constraints:} At extreme scale, costs for LLM API calls, database queries, and computation may become prohibitive. Cost optimization may be needed.
    \item \textbf{LLM Costs:} While the three-tier fallback architecture reduces LLM costs to approximately \$15 over the project duration---a 99.5\% reduction from the originally planned GPT-4---the Gemini API still incurs costs at scale.
    \item \textbf{Learner Motivation and Self-Discipline:} The system can structure learning and provide adaptive feedback, but it cannot directly motivate learners or prevent procrastination.
    \item \textbf{Hands-On Verification:} The system recommends resources and provides quizzes but does not directly assess learners' ability to write functional code.
\end{enumerate}

\section{Methodology and Approach}
\label{sec:methodology}

\subsection{Development Methodology: Agile-Scrum}

\projectname{} is developed using Agile-Scrum methodology, emphasising iterative delivery, continuous stakeholder feedback, and adaptive planning. Key practices include:
\begin{itemize}
    \item \textbf{Sprint Duration:} 2-week sprints, with each sprint focused on delivering incremental, testable features.
    \item \textbf{Daily Standups:} 15-minute daily synchronization meetings where team members report progress, identify blockers, and coordinate efforts.
    \item \textbf{Sprint Planning:} At sprint start, the team selects user stories from the product backlog, estimates effort, and commits to sprint goals.
    \item \textbf{Sprint Review:} At sprint end, completed features are demonstrated to stakeholders for feedback.
    \item \textbf{Sprint Retrospective:} The team reflects on process improvements, identifies what went well, and what needs adjustment.
\end{itemize}

\subsection{System Architecture: Microservices with Layered Design}

\projectname{} is architected as a modular microservices system with clear separation of concerns:

\begin{enumerate}
    \item \textbf{Frontend Layer (Presentation):} Web application (React 18 + TypeScript + Vite) with Zustand state management, Tailwind CSS, responsive design, and interactive DAG roadmap visualisation via React Flow. Flutter mobile application (Dart) for iOS/Android as a complementary access channel.
    \item \textbf{API Gateway Layer (Routing and Authentication):} Node.js/Express service handling HTTP request routing, JWT issuance and verification, Google/GitHub OAuth2 flows, rate limiting, and service-to-service proxying.
    \item \textbf{Learning Service (Core Business Logic):} Node.js/Express/Prisma service encapsulating all domain logic: enrollment management, learner progress tracking, Gatekeeper assessment evaluation, Mastery Decay computation, multi-path branching, certificate issuance, notifications, achievements, and instructor/admin operations.
    \item \textbf{AI Service (Tutor Model):} Dedicated Node.js service orchestrating the three-tier LLM provider chain (Phi-4 $\to$ Ollama $\to$ Gemini) with circuit breakers, Redis caching, SSE streaming, and structured prompt engineering for quiz generation, explanation generation, and AI Instructor chat.
    \item \textbf{Data Layer (Persistence):} PostgreSQL with Prisma ORM (learner profiles, ontology, progress, quiz attempts, certificates, achievements), Redis (AI response cache, session data).
    \item \textbf{Infrastructure Layer:} Docker Compose orchestration, Nginx reverse proxy (routing, SSE buffering control, SSL termination), and GitHub Actions CI/CD pipeline.
\end{enumerate}

\subsection{Hybrid AI Strategy: Skeleton \& Flesh}

The ontology and AI components implement the Skeleton \& Flesh architecture to solve hallucination:

\textbf{The Skeleton (Teacher Model):} A powerful LLM (Gemini 2.5 Flash or similar) generates an initial domain ontology---a directed acyclic graph where nodes represent learning concepts, edges represent prerequisite relationships, and metadata includes learning outcomes, estimated effort, and recommended resources. The generated ontology undergoes rigorous human expert review: domain experts verify each node for accuracy, completeness, and logical sequencing. Once verified, the skeleton is locked and immutable, stored in the database and served to the Tutor Model at runtime.

\textbf{The Flesh (Runtime, Tutor Model):} At runtime, when a learner encounters a node and requests an explanation, the Tutor Model (a lightweight LLM) generates an explanation grounded in the verified skeleton. It cannot invent new nodes or deviate from the skeleton structure. The Tutor Model also generates quiz questions targeting specific learning outcomes of that node, draws from relevant resources, and provides personalised elaborations based on learner context. By restricting generation to the scope of the skeleton, the Tutor Model prevents hallucination while maintaining flexibility and personalisation.

\subsection{Gatekeeper Pattern: Three-Tier Progression Logic}

Learners progress through the roadmap by passing performance thresholds defined by the Gatekeeper Pattern:

\begin{itemize}
    \item \textbf{Tier 1 -- Strong Mastery ($\ge 80\%$):} Quiz score $\ge 80\%$: Node marked \enquote{Done}, learner unlocks the next node. The learner receives a challenge project recommendation (a small real-world problem requiring application of the node's concepts). Node marked as \enquote{Mastered} (green status) with a 14-day decay timer initiated.
    \item \textbf{Tier 2 -- Marginal Mastery ($70$--$79\%$):} Quiz score 70--79\%: Node marked \enquote{Done}, learner unlocks next node, but node is flagged for spaced repetition review. Learners are not offered a challenge project immediately. Node marked as \enquote{Mastered} with a 7-day decay timer (faster decay due to lower initial mastery).
    \item \textbf{Tier 3 -- Insufficient Mastery ($< 70\%$):} Node remains locked. Score 50--69\% triggers a \texttt{resource\_swap} adaptation event (Tutor Model regenerates in tutorial-oriented style). Score $<50\%$ triggers prerequisite review. Score $<30\%$ escalates to the instructor via a flagged event.
\end{itemize}

\subsection{Spaced Repetition Integration: Mastery Decay Algorithm}

The Mastery Decay algorithm integrates Ebbinghaus's forgetting curve into the roadmap UI, making knowledge decay visible:

\textbf{Decay States:}
\begin{itemize}
    \item \textbf{Green (0--14 days):} \enquote{Mastered.} Learners are not prompted to review.
    \item \textbf{Yellow (14--30 days):} \enquote{Review Needed.} A micro-quiz (2--3 questions) is generated. Learners are notified.
    \item \textbf{Red ($>30$ days):} \enquote{Relearn.} A full-length quiz is triggered. Learners are strongly encouraged to re-engage.
\end{itemize}

\textbf{Micro-Quiz Logic:} When a node transitions from green to yellow or red, a micro-quiz is auto-generated targeting 2--3 core learning outcomes. Question formats include multiple choice, short answer, and code completion. If passed ($\ge 80\%$), the node timer resets for 14 days. If failed, the node is re-locked and resource adaptation is triggered.

\textbf{Visualization:} The learner's roadmap DAG is color-coded: green (mastered, current), yellow (review needed), red (relearn), gray (not yet attempted), blue (prerequisites not met). A timeline view shows decay over time.

\subsection{Multi-Path Branching Logic}

After completing foundational nodes (e.g., \enquote{Core Programming Fundamentals}), learners encounter a branching point:

\textbf{Branching Question Example:} \enquote{Which area aligns with your goals and interests? A) Frontend Development, B) Backend Development, C) Data Science \& AI}

Each path maintains prerequisite integrity. Paths converge at \enquote{Full-Stack Integration} (Frontend/Backend) or \enquote{AI/ML Systems Design} (Data Science) nodes, ensuring foundational coherence while respecting learner agency.

\subsection{Resource Generation: Ontology-Grounded AI Content}

Learning content is delivered through a two-layer approach that prevents hallucination by anchoring all AI output within the verified ontology:

\textbf{Offline Curation Layer:} During ontology authoring, domain experts attach curated reference links directly to each learning node. These resources are reviewed for accuracy and availability before publication and are served immediately to learners without AI mediation.

\textbf{Runtime AI Generation Layer:} When a learner enters a node, the Learning Service invokes the AI Service to generate a node-specific explanation. The prompt includes the node's title, learning outcomes, difficulty level, and prerequisite chain---all drawn from the verified skeleton. The Tutor Model may not introduce concepts or topics outside this context, preventing factual deviation. Explanations are streamed token-by-token to the client via SSE for immediate display.

\textbf{AI Instructor (Conversational):} Each learning node hosts a persistent AI Instructor chatbot. Learner questions are sent with the same node context block, ensuring answers remain grounded. Chat responses also stream via SSE, providing a responsive conversational experience without waiting for full response generation.

\textbf{Instructor Oversight:} Domain expert instructors can monitor flagged quiz events and learner analytics to identify nodes where learners consistently struggle, informing future ontology refinements or resource additions.

\subsection{Technology Stack and Justifications}

Table~\ref{tab:ch01_tech_stack} details the technology stack and rationale for each choice.

\begin{table}[htbp]
    \centering
    \caption{Technology stack and justifications}
    \label{tab:ch01_tech_stack}
    \begin{tabular}{lll}
        \toprule
        Component & Technology & Justification \\
        \midrule
        Frontend Web & React 18 + TypeScript + Vite & Type safety, fast HMR, large ecosystem \\
        Frontend Mobile & Flutter (Dart) & Cross-platform iOS/Android \\
        DAG Visualization & React Flow & Declarative node-graph with React \\
        API Gateway & Node.js + Express & Routing, JWT auth, rate limiting \\
        Learning Service & Node.js + Express + Prisma & Business logic, ORM type safety \\
        AI Service & Node.js + Express & AI provider orchestration \\
        Authentication & JWT + Google/GitHub OAuth2 & Stateless, social login support \\
        Database & PostgreSQL + Prisma ORM & ACID compliance, type-safe queries \\
        Cache & Redis & In-memory cache, sub-ms latency \\
        LLM --- Primary & Phi-4-Multimodal & Zero-cost, MIT licence \\
        LLM --- Fallback & Ollama (Qwen2.5-3B) & Local CPU inference, offline capable \\
        LLM --- Last Resort & Gemini 2.5 Flash API & High quality, cost-optimized \\
        AI Streaming & Server-Sent Events (SSE) & Real-time token streaming \\
        Reverse Proxy & Nginx & Load balancing, SSE buffering control \\
        Containerization & Docker + Docker Compose & Reproducible, portable deployment \\
        CI/CD & GitHub Actions & Automated testing and deployment \\
        Code Repository & GitHub & Open-source collaboration \\
        \bottomrule
    \end{tabular}
\end{table}

\subsection{LLM Selection Rationale}

The LLM architecture employs a three-tier fallback system designed to optimise for cost, capability, reliability, and suitability for the Ethiopian educational context.

\textbf{Tier 1 --- Primary: Phi-4-Multimodal (Microsoft).} Phi-4-Multimodal (5.6B parameters, $\sim$11GB in FP16) runs on Kaggle's free T4 GPU (16GB VRAM), providing LLM inference at zero financial cost---no cloud API fees, no credit card required, no usage limits. It is released under the permissive MIT license, ensuring no vendor lock-in. With native multimodal support (text, image, audio), it opens future possibilities for voice-based tutoring and diagram analysis. At $\sim$74\% MMLU, it provides adequate quality for quiz generation and explanation tasks. The primary limitation is inference speed ($\sim$16 tok/s) on Kaggle's T4 GPU, mitigated by response caching.

\textbf{Tier 2 --- Secondary: Ollama (Qwen2.5-3B, Alibaba).} The Ollama runtime serves Qwen2.5-3B on any modern CPU with no GPU required. At only 3 billion parameters ($\sim$6GB in FP16), inference completes in 30--90 seconds---slow but functional for cached or low-traffic scenarios. Despite its small size, Qwen2.5-3B achieves 65.9\% on MATH, 86.7\% on GSM8K, and 74.4\% on HumanEval. Crucially, Ollama provides complete offline capability---the system remains functional even when both the Phi-4 endpoint and the Gemini API are unavailable, a realistic scenario for Ethiopian deployments in areas with intermittent internet connectivity. The Ollama server is deployed alongside the AI Service container within the same Docker Compose stack.

\textbf{Tier 3 --- Tertiary: Gemini 2.5 Flash (Google).} Gemini 2.5 Flash is the last-resort fallback, used only when both local models are unavailable. With 88.4\% on MMLU and 82.8\% on GPQA Diamond, it outperforms both local alternatives by a wide margin. At 186--232 output tokens per second with a time-to-first-token of $\sim$0.44 seconds, it is 10--14x faster than Phi-4-Multimodal. The 1,048,576 token context window enables loading entire textbooks and comprehensive documentation into a single prompt. At \$0.15 per million input tokens and \$0.60 per million output tokens (with cached input at \$0.03/1M), it represents a 99.9\% cost reduction from the original GPT-4 budget estimate of \$450/month. The critical limitation is that it is proprietary and API-only, requiring internet connectivity.

\textbf{Cost Comparison with Original GPT-4 Plan:} Table~\ref{tab:ch01_llm_cost} compares the costs of the three-tier approach against the originally planned GPT-4 model.

\begin{table}[htbp]
    \centering
    \caption{Cost comparison: three-tier LLM architecture vs.\ original GPT-4 plan}
    \label{tab:ch01_llm_cost}
    \begin{tabular}{lrrrr}
        \toprule
        Model & Input Cost & Output Cost & Monthly Cost & 6-Month Cost \\
        \midrule
        GPT-4 (original plan) & \$30.00/1M & \$60.00/1M & $\sim$\$450 & \$2,700 \\
        Gemini 2.5 Flash & \$0.15/1M & \$0.60/1M & $\sim$\$1.50--3.00 & $\sim$\$15 \\
        Phi-4-Multimodal & \$0.00 & \$0.00 & \$0.00 & \$0.00 \\
        Qwen2.5-3B (Ollama) & \$0.00 & \$0.00 & \$0.00 & \$0.00 \\
        \bottomrule
    \end{tabular}
\end{table}

The three-tier architecture achieves a 99.5\% reduction in LLM operating costs compared to the original GPT-4 plan (\$2,700 $\rightarrow$ $\sim$\$15), while providing superior reliability through graceful degradation and complete offline capability.

\subsection{Development Phases}

The project is divided into six sequential phases over six months:

\begin{itemize}
    \item \textbf{Phase 1 (Weeks 1--4): Foundation \& Architecture.} System architecture documentation, database schema design, API specification, frontend component library setup, CI/CD pipeline configuration. Milestone: Architecture review approved, development environment setup complete.
    \item \textbf{Phase 2 (Weeks 5--8): Ontology Design \& Teacher Model.} Domain ontology for Frontend Development (30--50 learning nodes), ontology validation framework, Teacher Model implementation, expert verification process. Milestone: Pilot ontology created and verified.
    \item \textbf{Phase 3 (Weeks 9--12): Core Features --- Gatekeeper \& Tutor Model.} Gatekeeper Pattern implementation, Tutor Model (runtime LLM integration), quiz generation and evaluation logic, SERP API integration for supplementary resource discovery. Milestone: Gatekeeper logic tested, first resource recommendations surfaced.
    \item \textbf{Phase 4 (Weeks 13--16): Adaptation \& Persistence.} Mastery Decay algorithm, multi-path branching UI, learner progress persistence, spaced repetition notifications, learner dashboard. Milestone: Branching paths selectable in UI, progress data persisted.
    \item \textbf{Phase 5 (Weeks 17--20): Platforms \& Optimization.} Web interface fully functional, mobile interface MVP, performance optimization, scalability testing, cost analysis. Milestone: Web platform user testing with 10 learners, performance benchmarks meeting SLAs.
    \item \textbf{Phase 6 (Weeks 21--24): Validation \& Deployment.} Empirical evaluation with 50--100 Ethiopian learners, final bug fixes, production deployment, documentation and developer guides. Milestone: Pre/post assessments completed, statistical analysis of learning gains.
\end{itemize}

\section{Plan of Activities}
\label{sec:plan}

The project spans six months (24 weeks) divided into six phases with specific milestones, deliverables, and team responsibilities as detailed in Table~\ref{tab:ch01_plan}.

Table~\ref{tab:ch01_plan} provides a detailed weekly breakdown across 12 sprints.

\begin{longtable}{llllll}
    \caption{Detailed weekly breakdown (24 weeks)}
    \label{tab:ch01_plan}\\
    \toprule
    Weeks & Sprint & Focus Areas & Deliverables & Sprint Goals & Team Leads \\
    \midrule
    \endfirsthead
    \toprule
    Weeks & Sprint & Focus Areas & Deliverables & Sprint Goals & Team Leads \\
    \midrule
    \endhead
    \bottomrule
    \endfoot
    1--2   & Sprint 1  & Project kickoff, architecture design & Architecture doc v0.1, tech stack matrix, GitHub repo setup & Architecture review approved & Tech Lead \\
    3--4   & Sprint 2  & DB schema, API spec, CI/CD setup & PostgreSQL schema, OpenAPI spec, GitHub Actions, Docker setup & DB schema reviewed, build pipeline operational & Backend Lead \\
    5--6   & Sprint 3  & Ontology framework, Teacher Model prototype & Ontology JSON schema, Teacher Model prototype, expert interview notes & Ontology structure finalised, sample skeleton generated & AI Engineer \\
    7--8   & Sprint 4  & Pilot domain ontology (Frontend), expert verification & Frontend ontology (40 nodes), verification checklist (50 items) & Pilot ontology created and verified & AI Engineer \\
    9--10  & Sprint 5  & Gatekeeper Pattern implementation, quiz service & Gatekeeper logic, quiz service API, LLM quiz generation & Three-tier outcomes validated & Backend Lead + AI \\
    11--12 & Sprint 6  & Tutor Model, resource discovery (SERP API) & Tutor Model service, SERP API wrapper, learner feedback survey & Tutor Model generates explanations, resources fetched & Backend Lead + AI \\
    13--14 & Sprint 7  & Mastery Decay, spaced repetition, dashboard UI & Decay algorithm, micro-quiz generation, dashboard wireframes & Decay states correct, micro-quizzes triggered & Full-Stack \\
    15--16 & Sprint 8  & Multi-path branching, prerequisite validation & Branching logic, prerequisite graph validation, React branching UI & Paths selectable, validation prevents invalid paths & Frontend + Backend \\
    17--18 & Sprint 9  & React web interface (DAG, quiz, dashboard) & React components (React Flow DAG viewer, quiz UI, dashboard, AI panels) & Web interface functional, DAG interactive, quiz flow smooth & Frontend Lead \\
    19--20 & Sprint 10 & Flutter mobile interface & Flutter widgets (roadmap, quiz, notifications, offline DB) & Mobile app installable, offline tested, notifications triggered & Mobile Lead \\
    21--22 & Sprint 11 & Performance optimization, scalability testing & Performance benchmarks (load $<$2s, API $<$500ms), load test report (1000 concurrent) & SLAs met, system handles 1000 concurrent users & Tech Lead + DevOps \\
    23--24 & Sprint 12 & Empirical evaluation, production deployment & User study results, deployment checklist, documentation & Pre/post assessments analysed, production deployed & QA + Tech Lead \\
\end{longtable}

\section{Budget Breakdown}
\label{sec:budget}

Table~\ref{tab:ch01_budget} presents the six-month budget for development, testing, and initial deployment for 50--100 concurrent test users.




\section{Significance of the Project}
\label{sec:significance}

\subsection{Academic Significance}

\projectname{} contributes to three areas of academic research and practice in computer science and educational technology:

\textbf{1. Hybrid AI Architecture for Curriculum Design:} Traditional ITS faced a knowledge engineering bottleneck---each domain required thousands of hours of expert time. Modern LLMs solve scalability but introduce hallucination. \projectname{} operationalises a novel Skeleton \& Flesh variant of RAG specifically for curriculum design, bridging decades of ITS research with contemporary AI. This architecture is novel in its application to self-directed learning and could inform future EdTech research on verified, scalable knowledge representation.

\textbf{2. Assessment-Driven Resource Adaptation:} The Gatekeeper Pattern and three-tier remediation operationalise research from cognitive load theory~\cite{sweller1988} and multimedia learning theory~\cite{clark2016, mayer2009}. By grounding assessment outcomes in adaptive resource selection rather than just scoring feedback, \projectname{} advances the \enquote{knowing how to teach} dimension of adaptive tutoring, extending beyond item selection (IRT-based) to meta-cognitive strategy selection.

\textbf{3. Decay-Aware Roadmap Evolution:} By integrating Ebbinghaus's forgetting curve and modern spaced repetition research directly into curriculum roadmap interfaces, \projectname{} makes learning persistence a first-class system feature rather than a learner-managed tool. This is novel in positioning decay as a dynamic roadmap state (Green/Yellow/Red), enabling visualisations and system-level reminders that enhance retention science's practical impact.

\subsection{Regional and Socioeconomic Significance (Ethiopian Context)}

\textbf{1. Educational Access and Democratization:} In Ethiopia, access to quality technical education is limited by geographic constraints, institutional capacity, and cost. \projectname{} is free, web-accessible, and locally developed, removing these barriers. By generating personalized learning content grounded in expert-verified curricula and adapting it to each learner's performance, the system enables thousands of Ethiopian students to achieve technical competence without expensive mentorship or bootcamps.

\textbf{2. Employment Readiness and Economic Mobility:} The software engineering job market in Ethiopia is growing rapidly, with demand far exceeding supply. \projectname's focus on contemporary programming domains and assessment-driven progression ensures learners develop verifiable, job-aligned competencies, enhancing employment prospects and reducing skill gaps for employers.

\textbf{3. Reduction of Brain Drain and Talent Retention:} High-quality, accessible technical education increases the competitiveness of local talent and improves earning potential domestically, reducing incentives to emigrate. Retention of technical talent strengthens Ethiopia's domestic tech ecosystem.

\textbf{4. Open-Source and Knowledge Sharing Model:} \projectname{} is positioned as an open-source project with source code publicly available on GitHub. This enables Ethiopian tech communities, universities, and NGOs to contribute, adapt, and deploy locally, building local capacity in EdTech development.

\subsection{Technical and Practical Significance}

\textbf{Scalability of Adaptive Learning:} By decoupling curriculum knowledge (the skeleton) from runtime tutoring (the flesh), \projectname{} demonstrates how adaptive learning can scale beyond manual curation. The system can extend to additional domains by generating new skeletons with expert verification, without redesigning the tutoring engine.

\textbf{Search-Augmented Resource Discovery:} The use of SERP API for supplementary resource discovery demonstrates a practical approach to surfacing high-quality, domain-relevant learning materials at scale, applicable to other contexts (medical education, legal education, language learning) where curated search is preferable to open-web retrieval.

\textbf{Open-Source EdTech Contribution:} \projectname{} contributes to the open-source EdTech ecosystem alongside platforms like Moodle and Open edX, enabling researchers worldwide to build upon the work.

\subsection{User and Stakeholder Significance}

\textbf{Student Learners (Primary):} Students gain structured, adaptive guidance, reducing time-to-competence, improving retention, and increasing confidence. Multi-Path Branching respects learner agency and interests, increasing intrinsic motivation.

\textbf{Educators and Mentors (Secondary):} Teachers gain visibility into learner progress through dashboards and analytics, identify struggling learners (flagged by Gatekeeper failures), and provide targeted support.

\textbf{Employers:} An increased pool of technically skilled, job-ready graduates benefits employers. The system can be integrated into corporate training programs.

\textbf{Open-Source and Developer Communities (Indirect):} As an open-source platform, \projectname{} invites contributions from the developer community---new domain ontologies, AI prompt improvements, UI components, and language adaptations. This creates a compounding knowledge base that benefits both the project and contributors seeking practical EdTech experience.

\section{Outline of the Study}
\label{sec:outline}

This document is structured as follows:

\begin{description}
    \item[Chapter 1: Introduction (Current Chapter)] Articulates the problem landscape (tutorial hell, absence of adaptive feedback, lack of mastery-based progression, cognitive overload, generic roadmaps, hallucination). Defines objectives for system design, implementation, and validation. Specifies scope and limitations, outlines methodology (Agile-Scrum, microservices, Skeleton \& Flesh, Gatekeeper, Mastery Decay, Multi-Path Branching, SERP API integration), details the plan of activities, presents the budget, and justifies academic, regional, technical, and stakeholder significance.
    \item[Chapter 2: Literature Review] Synthesises research on ITS evolution, adaptive learning, knowledge representation (ontologies and RAG), assessment-driven progression, spaced repetition, personalisation, resource aggregation, and content curation. Analyses gaps in existing solutions and establishes three core lessons: ontology-guided LLM solves hallucination, assessment-driven resource adaptation extends tutoring, and decay-aware adaptation operationalises spacing research.
    \item[Chapter 3: Problem Analysis and Modeling] Provides detailed requirements elicitation methodology, specifies functional and non-functional requirements, and presents system models: use case diagrams, UML class diagrams, sequence diagrams, activity diagrams, and state machine diagrams.
    \item[Chapter 4: System Design and Architecture] Presents the detailed system architecture including microservices layers, Skeleton \& Flesh hybrid AI, Gatekeeper Pattern, Mastery Decay, Multi-Path Branching, database schema, API specifications, and deployment architecture.
    \item[Chapter 5: Implementation and Development] Documents development process, technologies used, code organisation, integration of the three-tier LLM stack and SERP API, SSE streaming implementation, certificate system, testing strategy, and lessons learned.
    \item[Chapter 6: Evaluation and Validation] Specifies evaluation methodology with empirical user study, defines success metrics (learning gains, engagement, usability), presents results, and discusses findings and limitations.
    \item[Chapter 7: Conclusions and Recommendations] Summarises project achievements, discusses how objectives were met, identifies technical and research contributions, and recommends future work and sustainability strategies.
\end{description}
